% from ega2-1-fr.tex
\oldpage[II]{9}
Comme on peut supposer $X$ associé à une $\mathcal{O}_S$-Algèbre en vertu de
(\hyperref[II.1.2.6-fr]{1.2.6}) et (\hyperref[II.1.3.1-fr]{1.3.1}), le
corollaire résulte de la construction de $X$ décrite dans
(\hyperref[II.1.3.1-fr]{1.3.1}).

\begin{example}[1.3.3]
\label{II.1.3.3-fr}
Soit $S$ le plan affine sur un corps $K$, où le point $0$ a été dédoublé
(\hyperref[I.5.5.11-fr]{I, 5.5.11})~; avec les notations de
(\hyperref[I.5.5.11-fr]{I, 5.5.11}), $S$ est réunion de deux ouverts affines
$Y_1$, $Y_2$~; si $f$ est l'immersion ouverte $Y_1\to S$,
$f^{-1}(Y_2)=Y_1\cap Y_2$ n'est pas un ouvert affine dans $Y_1$
(\emph{loc. cit.}), donc on a un exemple d'un schéma affine qui n'est pas
affine au-dessus de $S$.
\end{example}

\begin{corollary}[1.3.4]
\label{II.1.3.4-fr}
Soit $S$ un schéma affine~; pour qu'un $S$-préschéma $X$ soit affine au-dessus
de $S$, il faut et il suffit que $X$ soit un schéma affine.
\end{corollary}

\begin{corollary}[1.3.5]
\label{II.1.3.5-fr}
Soient $X$ un préschéma affine au-dessus d'un préschéma $S$, $Y$ un
$X$-préschéma. Pour que $Y$ soit affine au-dessus de $S$, il faut et il
suffit que $Y$ soit affine au-dessus de $X$.
\end{corollary}

On est aussitôt ramené au cas où $S$ est un schéma affine, et il en est alors
de même de $X$ (\hyperref[II.1.3.4-fr]{1.3.4})~; les deux conditions de
l'énoncé expriment alors que $Y$ est un schéma affine
(\hyperref[II.1.3.4-fr]{1.3.4}).

\begin{env}[1.3.6]
\label{II.1.3.6-fr}
Soit $X$ un préschéma affine au-dessus de $S$. Pour définir un préschéma $Y$
affine \emph{au-dessus de $X$}, il revient au même, en vertu de
(\hyperref[II.1.3.5-fr]{1.3.5}), de se donner un préschéma $Y$ affine
\emph{au-dessus de $S$}, et un $S$-morphisme $g:Y\to X$~; en d'autres termes
(\hyperref[II.1.3.1-fr]{1.3.1} et
\hyperref[II.1.2.7-fr]{1.2.7}), cela revient à se donner une
$\mathcal{O}_S$-Algèbre quasi-cohérente $\mathcal{B}$ et un homomorphisme
$\mathcal{A}(X)\to\mathcal{B}$ de $\mathcal{O}_S$-Algèbres (que l'on peut
envisager comme définissant sur $\mathcal{B}$ une structure de
$\mathcal{A}(X)$-Algèbre). Si $f:X\to S$ est le morphisme structural, on a
alors $\mathcal{B}=f_*(g_*(\mathcal{O}_Y))$.
\end{env}

\begin{corollary}[1.3.7]
\label{II.1.3.7-fr}
Soit $X$ un préschéma affine au-dessus de $S$~; pour que $X$ soit de type fini
sur $S$, il faut et il suffit que la $\mathcal{O}_S$-Algèbre quasi-cohérente
$\mathcal{A}(X)$ soit de type fini
(\hyperref[I.9.6.2-fr]{I, 9.6.2}).
\end{corollary}

Par définition (\hyperref[I.9.6.2-fr]{I, 9.6.2}), on est ramené au cas où
$S$ est affine~; alors $X$ est un schéma affine
(\hyperref[II.1.3.4-fr]{1.3.4}), et si $S=\operatorname{Spec}(A)$,
$X=\operatorname{Spec}(B)$, $\mathcal{A}(X)$ est la
$\mathcal{O}_S$-Algèbre $\widetilde{B}$~; comme
$\Gamma(U,\widetilde{B})=B$, le corollaire résulte de
(\hyperref[I.9.6.2-fr]{I, 9.6.2}) et
(\hyperref[I.6.3.3-fr]{I, 6.3.3}).

\begin{corollary}[1.3.8]
\label{II.1.3.8-fr}
Soit $X$ un préschéma affine au-dessus de $S$~; pour que $X$ soit réduit il
faut et il suffit que la $\mathcal{O}_S$-Algèbre quasi-cohérente
$\mathcal{A}(X)$ soit réduite
(\hyperref[0.4.1.4-fr]{0, 4.1.4}).
\end{corollary}

En effet, la question est évidemment locale sur $S$~; en vertu de
(\hyperref[II.1.3.2-fr]{1.3.2}), le corollaire résulte de
(\hyperref[I.5.1.1-fr]{I, 5.1.1}) et
(\hyperref[I.5.1.4-fr]{I, 5.1.4}).

\subsection{Faisceaux quasi-cohérents sur un préschéma affine au-dessus de
$S$}
\label{subsection:II.1.4-fr}

\begin{proposition}[1.4.1]
\label{II.1.4.1-fr}
Soient $X$ un préschéma affine au-dessus de $S$, $Y$ un $S$-préschéma,
$\mathcal{F}$ (resp. $\mathcal{G}$) un $\mathcal{O}_X$-Module (resp. un
$\mathcal{O}_Y$-Module) quasi-cohérent. Alors l'application
$(h,u)\mapsto(\mathcal{A}(h),\mathcal{A}(u))$ de l'ensemble des morphismes
$(Y,\mathcal{G})\to(X,\mathcal{F})$ dans l'ensemble des di-homomorphismes
$(\mathcal{A}(X),\mathcal{A}(\mathcal{F}))\to
(\mathcal{A}(Y),\mathcal{A}(\mathcal{G}))$
(\hyperref[II.1.1.2-fr]{1.1.2} et
\hyperref[II.1.1.3-fr]{1.1.3}) est bijective.
\end{proposition}

La démonstration suit exactement la même marche que celle de
(\hyperref[II.1.2.7-fr]{1.2.7}) en utilisant
(\hyperref[I.2.2.5-fr]{I, 2.2.5}) et
(\hyperref[I.2.2.4-fr]{I, 2.2.4}), et les détails en sont laissés au
lecteur.

\begin{corollary}[1.4.2]
\label{II.1.4.2-fr}
Si, en outre des hypothèses de (\hyperref[II.1.4.1-fr]{1.4.1}), on suppose
$Y$ affine au-dessus de $S$, alors, pour que $(h,u)$ soit un isomorphisme,
il faut et il suffit que $(\mathcal{A}(h),\mathcal{A}(u))$ soit un
di-isomorphisme.
\end{corollary}

